\subsection{Wesentliche Herausforderungen und Risiken im Projekt}
\label{subsec:WesentlicheHerausforderungen}

\par Im Projektverlauf stellte insbesondere die Diskrepanz zwischen dem ursprünglich propagierten \gls{gfa} und der tatsächlichen Komplexität der Altsystem-Integration eine zentrale Herausforderung dar. Ich begleitete maßgeblich den strategischen Schwenk hin zu einer realistischeren Auslieferungsplanung im Rahmen eines \gls{bfa} gegenüber der Außenwelt und steuerte die daraus resultierenden Scope-Anpassungen durch konsequentes Risikomanagement. Eine weitere kritische Hürde war der Mangel an historischem \gls{e2e}-Prozesswissen im \gls{fb} und fehlendem internen \gls{s4}-Konw-How in der \gls{s4}, dem ich durch den gezielten Aufbau spezialisierter Ressourcen und Wissensformate entgegenwirkte. In Krisensituationen, etwa bei qualitativen Defiziten externer Partner, sicherte ich durch moderierte Klärungsrunden und eine transparente Change-Kommunikation die Stabilität des Gesamtprogramms. Die aktive Gestaltung des Wandels hin zu einer \gls{pcr}-Architektur war dabei entscheidend, um die Akzeptanz der neuen Prozesse bei den betroffenen Stakeholdern zu festigen und Konflikte frühzeitig zu neutralisieren.

\begin{comment}
\par Im Verlauf des Programms \ac{r2h} traten zahlreiche fachliche, technische und organisatorische Herausforderungen auf. Früh zeigte sich, dass die initiale Erhebungsphase die tatsächliche Komplexität nicht vollständig erfasst hatte. Besonders deutlich wurde dies an den konkurrierenden Prämissen zwischen dem fachlichen \glqq Greenfield-Ansatz\grqq{} und der Erwartung einer weitgehend unveränderten Integration in die bestehende Altsystemwelt. Dies führte zu Verschiebungen in der Auslieferungsplanung, etwa durch die Aufspaltung der ursprünglichen Welle 1 und die Einführung neuer SAP-Komponenten wie \ac{mbc} und \ac{aif}. Die Steuerung dieser Abweichungen erfolgte durch konsequente Plananpassungen, enge Abstimmungen mit den fachlichen \ac{wsl} und dem \ac{ws} Migration sowie ein kontinuierliches Status- und Risikomanagement.

\par Eine weitere zentrale Herausforderung war die Identifikation und Allokation geeigneter Ressourcen. Da die Implementierung des \ac{r3} über 25 Jahre zurücklag, fehlte erfolgskritisches Know-how in der Organisation. Daher mussten benötigte Skills für die Umsetzung der \ac{ricefw}-Objekte fortlaufend bewertet und aufgebaut werden. Dies geschah durch die Analyse technischer Anforderungen wie \ac{aif} oder \ac{btp}, die Bewertung interner sowie externer Kompetenzen und den gezielten Einsatz spezialisierter Ressourcen. Ergänzt wurde dies durch Skill-Mapping sowie die enge Zusammenarbeit mit Lieferanten, Fachbereichen, den internen \ac{ais}-Mitarbeitenden und der Programmleitung.

\par Zur Beherrschung der Komplexität wurde ein umfassendes Risikomanagement etabliert, das Risiken kontinuierlich identifizierte, bewertete und priorisierte. Wesentliche Maßnahmen umfassten den strategischen Schwenk von einem \glqq Greenfield\grqq{} hin zu einem \glqq Brownfield-Ansatz\grqq{} zur Reduzierung technischer Risiken und die frühzeitige Eskalation in die zuständigen Gremien. Durch das abgestimmte Zusammenspiel von Risiko-, Abhängigkeits- und Testmanagement konnten Risiken entlang der gesamten Prozess- und Systemkette betrachtet und wirksam gesteuert werden. Regelmäßige Steering-Meetings dienten der Validierung der getroffenen Maßnahmen.

\par Das Changemanagement gewann durch Scope-Anpassungen, technische Änderungen und strategische Entscheidungen zunehmend an Bedeutung. Die Einführung zusätzlicher SAP-Module, die Integration der \ac{adn} und die Umstellung auf eine \glqq Public Cloud Ready\grqq{}-Zielarchitektur erforderten fortlaufende Anpassungen von Prozessen und Arbeitsweisen. Als \ac{wsl} lag es in meiner Verantwortung, diese Veränderungen transparent zu kommunizieren, deren Auswirkungen zu bewerten und die Mitarbeitenden zu begleiten. Durch klare Kommunikationsstrukturen und regelmäßige Abstimmungen konnten Unsicherheiten reduziert und Akzeptanz für das \ac{s4} geschaffen werden.

\par Konflikte und Krisen entstanden insbesondere im Zusammenspiel mit den in den ersten Wellen betroffenen Gesellschaften \ac{ais}, \ac{avm} und \ac{adn}, zwischen fachlichen Workstreams hinsichtlich Lieferprioritäten sowie bei Verzögerungen durch den externen Implementierungspartner. Diese äußerten sich etwa in unsynchronen Zeitplänen, unzureichenden Zulieferungen oder divergierenden Qualitätsmaßstäben. Als Workstream Lead war es notwendig, Konflikte frühzeitig zu erkennen, anzusprechen und geeignete Maßnahmen einzuleiten. Dazu gehörten moderierte Klärungsrunden, die Abstimmung realistischer Zeitpläne, die Eskalation wiederkehrender Probleme sowie die konsequente Verfolgung der Lieferverpflichtungen. Dadurch konnten kritische Situationen stabilisiert und negative Auswirkungen auf das Gesamtprogramm verringert werden.
\end{comment}